\documentclass[12pt]{article}
\usepackage{amsmath,amssymb,amsthm}
\usepackage[margin=1in]{geometry}
\usepackage{booktabs}
\usepackage{hyperref}

\newtheorem{theorem}{Theorem}[section]
\newtheorem{corollary}[theorem]{Corollary}
\newtheorem{proposition}[theorem]{Proposition}
\newtheorem{definition}[theorem]{Definition}
\theoremstyle{remark}
\newtheorem{remark}[theorem]{Remark}

\newcommand{\CC}{\mathbb{C}}
\newcommand{\RR}{\mathbb{R}}
\newcommand{\QQ}{\mathbb{Q}}
\newcommand{\Lean}{\textsc{Lean\,4}}

\title{P $\ne$ NP as Abel--Ruffini:\\
  The Algebraic Solve/Check Separation\\
  at Dimension 5}
\author{Mingu Jeong\\Independent Researcher
\and Claude\\Anthropic}
\date{\today}

\begin{document}
\maketitle

\begin{abstract}
We identify the P $\ne$ NP problem as the computational
analog of the Abel--Ruffini theorem (1824). \textbf{Solve}
(finding individual roots by radicals) is impossible for
degree $\ge 5$ because $S_5$ is non-solvable.
\textbf{Check} (computing symmetric functions via Vieta)
always works at any degree. The separation Solve $\ne$ Check
at $d = 5$ is proven; the computational P $\ne$ NP remains
open because $\QQ$ admits numerical approximation (Newton's
method), creating a gap absent in $\CC$. We formalize this
in \Lean{} and classify the problem as
standard $(0, 4)$, algebraic $(1, 2)$ \cite{Paper9}.
\end{abstract}

\section{Introduction}

The P $\ne$ NP conjecture asks whether every problem whose
solution can be \emph{verified} in polynomial time can also
be \emph{solved} in polynomial time. We restate this in
algebraic terms.

\section{Solve vs Check}

\begin{definition}
\textbf{Solve}$(P)$: express individual roots of a polynomial $P$
in terms of radicals ($+, -, \times, \div, \sqrt[n]{\phantom{x}}$).

\textbf{Check}$(P)$: compute symmetric functions of the roots
(sums, products, traces) via Vieta's formulas.
\end{definition}

\begin{theorem}[Abel--Ruffini, 1824]
For $\deg P \le 4$: Solve is possible ($S_n$ solvable).
For $\deg P \ge 5$: Solve is impossible ($S_5$ non-solvable).
Check is always possible (Vieta, any degree).
\end{theorem}

\begin{corollary}[Algebraic P $\ne$ NP]
There exist problems (degree $\ge 5$ polynomials) where
Check works but Solve does not. The boundary is $d = 5$,
the same dimension as DRLT.
\end{corollary}

\section{The Galois Obstruction}

$S_5$ is the first non-solvable symmetric group.
The obstruction is $A_5$, the alternating group:
$|A_5| = 60 = 2^2 \times 3 \times 5$,
built from exactly the DRLT atoms $\{2, 3, 5\}$.

\begin{theorem}[Galois--DRLT correspondence]
Solvable $\Leftrightarrow$ physically incomplete.
Non-solvable $\Leftrightarrow$ physically complete.
The boundary $d = 5$ is simultaneously the solvability
threshold and the completeness threshold.
\end{theorem}

\section{Why Computational P $\ne$ NP Remains Open}

In $\CC$: ``unsolvable = uncomputable'' (no gap).
Abel--Ruffini closes the question.

In $\QQ$: ``unsolvable $\ne$ uncomputable'' (gap exists).
Newton's method finds roots to arbitrary precision
even for degree $\ge 5$. The \emph{numerical
approximation gap} prevents Abel--Ruffini from
directly implying computational P $\ne$ NP.

\begin{remark}
This gap is the difference between $\CC$ (algebraically
closed, no approximation) and $\QQ$ (not algebraically
closed, approximation works). The P $\ne$ NP problem
is hard because $\QQ$ has this escape route.
\end{remark}

\section{Physics Needs Only Check}

All physical observables are symmetric functions of the
Gram matrix:
$\mathrm{Tr}(G^k)$, $\det(G_h)$, $|G_{ij}|^2$.
No individual eigenvalue appears in any physical quantity.
Physics is immune to the Abel--Ruffini barrier.

\section{Spectral Complexity}

Standard (computational): $(h, l) = (0, 4)$.
Algebraic: $(h, l) = (1, 2)$ (proven, Abel--Ruffini).
The algebraic version has been solved since 1824 \cite{Paper9}.

\section{Machine Verification}

\texttt{SolveCheck.lean}: \texttt{algebraic\_P\_ne\_NP},
\texttt{boundary\_is\_5}, \texttt{d5\_triple},
\texttt{physics\_immune\_to\_abel\_ruffini}.
Zero \texttt{sorry}.

\begin{thebibliography}{9}
\bibitem{Paper9} M.~Jeong and Claude, ``Spectral Complexity,''
  DRLT Paper~9 (2026).
\bibitem{Abel} N.~H.~Abel, ``M\'emoire sur les \'equations
  alg\'ebriques,'' 1824.
\end{thebibliography}
\end{document}